\startfirstchapter{Introduction}
\label{chapter:introduction}

Canada faces a growing legacy of abandoned oil wells and contaminated petrochemical sites. Surveys estimate that 15,682 abandoned wells require plugging and/or land reclamation in Alberta and Saskatchewan alone \cite{est_cost_of_cleaning_canada}. The presence of petroleum hydrocarbons (PHC) in soil poses serious risks to the environment, agricultural productivity, and human health in surrounding communities \cite{PHCBadForCrops,PHCBadForHumans,PHCBadForEnv}. As the world transitions away from fossil fuels, more facilities will be decommissioned under inadequate closure procedures, and the scale of this problem will only increase.

In situ remediation technologies artificially accelerate biodegradation within the soil without excavating the contaminated material \cite{remTechReview}. Biostimulation, one such proven approach, introduces nutrients directly into PHC-contaminated sites to promote microbial activity and hasten the remediation process. Compared to ex situ methods which require physically removing and treating soil off-site; in situ techniques are substantially less costly and less environmentally disruptive \cite{LioraTech}. However, they demand a much richer understanding of soil conditions to treat contaminated zones accurately, and they rely on in situ measurements to verify treatment progress. Historically, these measurements were collected by hand, yielding data that was sparse both spatially and temporally. In recent years, IoT sensor networks deployed at contaminated sites have enabled continuous monitoring of pollutant concentrations and environmental conditions \cite{LioraRemoteSensing,atmMeasInSoil}, providing dynamic estimates of how a PHC plume is distributed and evolving over time.

This adoption of continuous remote sensing has largely resolved the temporal sparsity problem. Spatial sparsity, however, remains a significant challenge: deploying and installing dense sensor arrays is expensive, leaving large gaps in coverage. Without an accurate spatial map of the PHC plume, remediators cannot effectively target treatment, resulting in wasted resources and land that remains unusable long after intervention.

Spatial interpolation offers a principled way to reconstruct plume geometry from sparse measurements, yet the wide variety of available methods---originating from fields as different as geostatistics, numerical weather prediction, and inverse modelling---makes systematic comparison difficult, particularly within a soil remediation context. This thesis addresses that gap by evaluating several interpolation approaches on both synthetic and real-world datasets to identify the most effective method for petrochemical remediation in Canada.

The methods compared are Inverse Distance Weighting (IDW), Kriging, 4D-Variational Data Assimilation (4D-Var), the Gauss--Levenberg--Marquardt algorithm (GLM), and the Iterative Ensemble Smoother (IES). These naturally divide into two categories. IDW and Kriging are distance-based interpolators that estimate concentration at unsampled locations as a weighted sum of nearby observations according to a chosen spatial model. By contrast, 4D-Var, GLM, and IES are parameter estimation algorithms that fit a physics-based transport model to temporal measurements; once calibrated, the model provides concentration estimates at any unobserved location within the domain.

Each method is evaluated on two datasets. A synthetic soil contaminant dataset with dense ground truth is used to tune hyperparameters and systematically assess robustness to measurement noise, soil heterogeneity, sensor dropout, and spatial and temporal sparsity. Building on this controlled evaluation, the methods are then applied to real sensor data collected in the Canadian prairies using leave-$n$-out cross-validation.

% TODO: Summarize key findings here, then provide a section-by-section roadmap.
% e.g., "The results indicate that ... Chapter 2 reviews related work, Chapter 3 describes ..."